\documentclass[11pt,a4paper]{article}
\usepackage[utf8]{inputenc}
\usepackage[margin=1in]{geometry}
\usepackage{xcolor}
\usepackage{listings}
\usepackage{enumitem}
\usepackage{longtable}
\usepackage{booktabs}
\usepackage[hidelinks]{hyperref}
\usepackage{titlesec}

\definecolor{codebg}{rgb}{0.96,0.96,0.96}
\lstset{
  basicstyle=\ttfamily\small,
  backgroundcolor=\color{codebg},
  breaklines=true,
  frame=single,
  columns=fullflexible,
  showstringspaces=false,
  xleftmargin=4pt, xrightmargin=4pt
}
\titleformat{\section}{\large\bfseries}{\thesection}{0.6em}{}
\setlength{\parskip}{0.5em}
\setlength{\parindent}{0pt}

\title{\textbf{Accomplish $\rightarrow$ SRIM (DigiBull) Web Migration}\\[2pt]
\large Technical Report: Conversion, Analysis, and Implementation}
\author{DigiBull.ai --- SRIM 3.2 Programme}
\date{June 2026}

\begin{document}
\maketitle

\begin{abstract}
This document records the migration of the \emph{Accomplish} desktop application
(Electron) to a browser-served \emph{local web} application, rebranded as
\textbf{SRIM} for DigiBull.ai. It describes the original architecture, the initial
browser conversion, the subsequent gap analysis and reverse-engineering of a
newer packaged release, and the features and fixes implemented to bring the web
build to a working state. It also documents the Model Context Protocol (MCP)
integration layer, the DigiBull UI update, task-aware Auto model routing, local
model download/runtime UX, and how to use the system. The account is factual and
neutral; named contributions are recorded only to attribute work.
\end{abstract}

\tableofcontents
\vspace{1em}

% ---------------------------------------------------------------
\section{Original Application (Electron)}
The upstream Accomplish app is a pnpm monorepo:

\begin{itemize}[leftmargin=1.4em,topsep=2pt]
  \item \texttt{apps/desktop} --- Electron shell: main process + preload bridge.
  \item \texttt{apps/web} --- standalone React UI (Vite + Zustand).
  \item \texttt{apps/daemon} --- background Node.js process (no Electron imports).
  \item \texttt{packages/agent-core} --- core logic, storage, MCP tooling, the
        OpenCode agent adapter.
\end{itemize}

\textbf{Data flow.} The React UI calls \verb|window.accomplish.*| (a bridge of
roughly 462 methods). In Electron this was served by the preload
\verb|contextBridge| over \verb|ipcRenderer|; the main process forwarded work to
the forked daemon, which runs the OpenCode agent runtime and persists state in
SQLite (\verb|better-sqlite3|). Native/desktop dependencies include OpenCode,
Playwright, Baileys (WhatsApp), \verb|electron-store|, and \verb|electron-updater|.

The renderer (\texttt{apps/web/src/client}) contains no Electron or Node-only
imports; the desktop coupling lived entirely in the main process, preload, and
daemon. This is what made a browser migration feasible.

% ---------------------------------------------------------------
\section{Baseline Artifact and Attribution}
The baseline for the converted web project is the archive:

\begin{lstlisting}
LOCAL/SRIM-UI/accomplish_Transformed_Chaithanya_Original.zip
\end{lstlisting}

This archive is the original Chaithanya transform artifact (approximately
759 MB) and is treated as the starting point for the current SRIM web working
tree. Subsequent analysis, gap remediation, UI changes, provider fixes, and
documentation updates continue from that baseline. This distinction is important
for fair attribution: the web conversion baseline is credited to Chaithanya; the
later SRIM verification, remediation, reverse-engineering notes, and current
DigiBull UI/model updates are credited separately.

% ---------------------------------------------------------------
\section{Phase 1 --- Conversion to a Local Web App}
\emph{Initial browser conversion (Chaithanya).} The Electron window/host was
removed so the UI can run in an ordinary browser tab against a local backend.
Key additions:

\begin{itemize}[leftmargin=1.4em,topsep=2pt]
  \item \texttt{apps/daemon/src/browser-api-server.ts} --- an HTTP + Server-Sent
        Events (SSE) server on \verb|127.0.0.1:9234|, letting a browser tab reach
        the daemon directly without Electron.
  \item \texttt{apps/web/src/client/lib/accomplish-backend.ts} --- a browser
        transport that recreates \verb|window.accomplish| as a \verb|Proxy| over
        HTTP (\verb|/rpc|) and SSE (\verb|/events|).
  \item \texttt{session.ts} + \texttt{LoginPage.tsx} --- a login screen using a
        NetBird Personal Access Token; DigiBull/SRIM branding.
  \item Daemon dev scripts (tsup build + run) and supporting store changes.
  \item Delivery baseline:
        \verb|LOCAL/SRIM-UI/accomplish_Transformed_Chaithanya_Original.zip|.
\end{itemize}

At the end of this phase the UI loaded in the browser, but only part of the
\verb|window.accomplish| surface was wired through the new transport; the
remainder defaulted to safe no-ops or stubs.

% ---------------------------------------------------------------
\section{Phase 2 --- Verification, Gap Analysis, Reverse-Engineering}
\emph{Analysis and remediation (Nagabhushana).} The converted build was compared
against the original repository to establish exactly what was and was not wired.

\subsection{Inventory and gap analysis}
Two analyses were produced (see Artifacts):
\begin{itemize}[leftmargin=1.4em,topsep=2pt]
  \item \textbf{migration\_inventory.md} --- maps every Electron/Node API surface
        and where it lived.
  \item \textbf{missing\_features\_gap\_analysis.md} --- audits the
        \verb|window.accomplish| contract ($\sim$148 capabilities) against what
        the browser path actually wires. Result: $\sim$96 wired, $\sim$52 missing,
        bucketed by remediation type:
  \begin{description}[leftmargin=1.4em,topsep=1pt]
    \item[A --- Trivial-wire (24)] backed by an existing daemon service; add a
          channel.
    \item[B --- Needs daemon port (15)] logic lived in Electron main.
    \item[C --- Browser-native (6)] file pickers, clipboard, downloads, capture.
    \item[D --- Not possible in browser (7)] tray, updater, native window control.
  \end{description}
\end{itemize}

\subsection{Reverse-engineering the 0.4.14 release}
A packaged installer (\verb|Accomplish-0.4.14-win-x64.exe|, $\sim$294\,MB) was
examined to see whether a newer build carried features absent from the source
snapshot. Method:

\begin{enumerate}[leftmargin=1.4em,topsep=2pt]
  \item Extract the NSIS installer with 7-Zip $\rightarrow$
        \verb|$PLUGINSDIR/app-64.7z|.
  \item Extract that archive $\rightarrow$ \verb|resources/app.asar| (35\,MB,
        4313 files).
  \item Parse the \verb|asar| header (a dependency-free Node script) and diff the
        IPC channel string literals in the packaged main bundles against the
        source tree.
\end{enumerate}

\textbf{Finding.} Although the installer's internal version is also \texttt{0.1.0},
its code is newer than the source snapshot: $\sim$124 channel literals appear in
the \texttt{.exe} but not in source. Verified feature clusters absent from source
include an embedded \emph{Browser Panel}, \emph{Triggered Tasks}, local
\emph{llama.cpp} inference, \emph{feature flags}, and auto model-routing.
Conclusion: these can only be ported from 0.4.14 \emph{source} (the installer
contains only minified bundles), and several are native/desktop features.
Artifacts of this analysis are in \verb|LOCAL/exe-inspect-0.4.14/| and
\verb|exe_0.4.14_feature_delta.md|.

% ---------------------------------------------------------------
\section{Phase 3 --- Implemented Features and Fixes}
The following were implemented on the converted build. Each was verified at the
TypeScript transpile level; full runtime verification requires running the app
under Node~24 (see \S\ref{sec:run}).

\subsection{Defect fixes}
\begin{longtable}{p{0.32\textwidth} p{0.60\textwidth}}
\toprule
\textbf{Area} & \textbf{Fix} \\
\midrule
\endhead
Agent task execution & The daemon never passed provider API keys to the OpenCode
process, so native-provider tasks (OpenAI/Anthropic/\dots) failed with no
credentials. The spawn environment now injects \verb|ANTHROPIC_API_KEY|,
\verb|OPENAI_API_KEY|, etc. via \verb|buildOpenCodeEnvironment|. \\
Provider connect crash & \verb|validateApiKeyForProvider| was unwired and
returned \verb|null|, so the UI crashed reading \verb|.valid|. Wired to
agent-core validation. \\
Key validation false-negatives & A retired probe model (e.g.
\texttt{claude-3-haiku}) caused valid keys to be rejected. Validation now treats
only \verb|401/403| as invalid (auth), not other statuses. \\
API-key quoting & Keys pasted with surrounding quotes broke auth headers; a
sanitizer strips them on save. \\
Model listing & \verb|fetchProviderModels| was forwarded with the wrong
signature; re-implemented to resolve endpoint + stored key like the original. \\
Hardcoded addresses & The browser-backend address was hardcoded in 5 places;
centralized into one registry (see below). \\
Silent failures & Failed tasks stored no reason; the failure is now persisted as
a system message and visible in history. \\
Auto routing stalls & Auto model routing is guarded by a timeout. If provider
settings or model switching hangs, task submission continues on the current
selection and a warning is logged, preventing the UI from staying on
\emph{Executing...}. \\
Provider status nulls & HuggingFace/local provider status is normalized before
rendering so a missing server-status response cannot crash the Providers tab. \\
\bottomrule
\end{longtable}

\subsection{New capabilities}
\begin{longtable}{p{0.32\textwidth} p{0.60\textwidth}}
\toprule
\textbf{Feature} & \textbf{Description} \\
\midrule
\endhead
Centralized address registry & \verb|config/endpoints.ts| (client) +
\verb|BROWSER_API_PORT|/\verb|WEB_DEV_ORIGIN| (agent-core), env-overridable;
production defaults to same-origin --- ``migration-proof'' addressing. \\
Dev test-login & A local login bypass (token or user/password) for testing
without the NetBird service; auto-disabled when \verb|NODE_ENV=production|. \\
Bucket-A wiring (24) & Settings config getters/setters, model fetchers,
validators, and secrets helpers wired through the browser backend. \\
Provider parity (bucket B) & Vertex credential save/get, Speech (ElevenLabs)
config/validate/transcribe, and OpenAI ``Sign in with ChatGPT'' OAuth. \\
Appearance / Customize & Accent palettes, dim (medium) mode, animation toggle,
background patterns; persisted client-side. \\
LOCAL LLM advisor & Detects CPU/RAM (daemon \verb|os|) + GPU/network (browser)
and recommends local model sizes; shown atop the providers panel. \\
DigiBull chat rebrand & Home prompt, placeholder, helper copy, and hover/motion
treatment were updated around the DigiBull identity without hardcoding all copy
in component logic. \\
Task-aware Auto model selector & The model selector exposes
\textbf{Auto $\cdot$ by task type}. At task-submit time the client classifies
the prompt (chat, coding, research, planning, analysis), prefers fast/cheap
models for simple chat, and stronger models for code, research, and multi-step
work. If routing cannot resolve, the existing selected model is used. \\
Inline AI question card & Question-type permission requests render as an inline
\textbf{DigiBull Question} card so the agent can ask for clarification and the
user can answer without a full-screen modal. File/tool permissions still use the
stronger modal. \\
Local model downloads & A local model download panel sits under the device
checker in Providers. It shows selected model state (not selected, ready to
download, downloading, installed, failed), download progress, and local server
running/idle state. \\
Runtime target selector & The HuggingFace Local flow exposes Auto, GPU, WebGPU,
and CPU runtime targets with helper text. Manual GPU index routing
(GPU~0/GPU~1) is not exposed yet; GPU modes use automatic GPU selection. \\
Performance toggle & \verb|ACCOMPLISH_BROWSER_MODE=none| skips the per-task
browser MCP/Chromium spin-up for faster chat startup. \\
Test-automation skill & A \verb|/test-automation| skill: spec $\rightarrow$ test
cases $\rightarrow$ run $\rightarrow$ compare $\rightarrow$ report. \\
Provenance & \verb|system:about| channel + visible login credit + \verb|AUTHORS.md|
(open, queryable authorship). \\
Developer mode + MCP manager & A gated Developer section with an MCP connection
manager (pre-seeded with the SRIM stack) and an opt-in integration API key. \\
\bottomrule
\end{longtable}

% ---------------------------------------------------------------
\section{Implementation Coverage}
\begin{longtable}{p{0.30\textwidth} p{0.16\textwidth} p{0.46\textwidth}}
\toprule
\textbf{Group} & \textbf{Status} & \textbf{Notes} \\
\midrule
\endhead
Bucket A (24) & Done & wired + transpile-verified \\
Bucket B portable & Partly done & Vertex creds, Speech, OpenAI OAuth, validation,
local model UI/runtime selection \\
Bucket B native & Open & manual GPU index selection, gcloud Vertex project APIs,
remaining native desktop-only services \\
OAuth (Copilot/Slack) & Open & need provider accounts to build/verify \\
Bucket C & Partly done & pickers/openExternal native; \verb|exportLogs|,
capture* open \\
Bucket D (7) & N/A & not possible in a browser tab (correctly stubbed) \\
0.4.14 features & Open & require 0.4.14 source (browser-panel, llama.cpp, etc.) \\
\bottomrule
\end{longtable}

% ---------------------------------------------------------------
\section{MCP Integration}\label{sec:mcp}
There are two directions. See \verb|MCP_GUIDE.md| for step-by-step usage.

\textbf{Outbound (Accomplish $\rightarrow$ MCP servers).} In
\emph{Settings $\rightarrow$ General $\rightarrow$ Developer}, enable Developer
mode and add remote MCP servers (name + URL). These are stored as connectors and
injected into the OpenCode config (\verb|buildMcpServers|), so the agent uses
them as tools. The manager is pre-seeded with quick-add templates for the SRIM
stack (Dify, FastMCP, MCP Jungle, Postgres, Apache Superset, and the tool MCPs).

\textbf{Inbound (external workflows $\rightarrow$ Accomplish).} A generated
integration API key lets external systems (e.g. Dify/FastMCP) call the local API:
\begin{lstlisting}
POST http://127.0.0.1:9234/rpc
Authorization: Bearer <integration-api-key>
Content-Type: application/json
{ "channel": "task:start", "args": [ { "prompt": "...", "files": [] } ] }
\end{lstlisting}
The key is honored \emph{only} when the daemon is started with
\verb|ACCOMPLISH_ENABLE_INTEGRATION_API=1| (opt-in), and the server binds to
\verb|127.0.0.1| only.

% ---------------------------------------------------------------
\section{DigiBull UI and Model UX Update}
This update focuses on making SRIM feel understandable while it works, especially
when the selected model is automatic or the agent needs one more answer from the
user.

\subsection{Chat entry and motion}
The home page prompt is now DigiBull-branded and the task input bar has a small
hover/focus lift, gradient edge, and status pill when Auto mode is active. The
motion is intentionally restrained: it is used to show focus and readiness, not
to animate every surface.

\subsection{Auto model routing}
The selector includes \textbf{Auto $\cdot$ by task type}. This is not only a
label: the client resolves an actual provider/model before \verb|task:start| by
reading connected provider settings and classifying the task prompt. Current
heuristics:

\begin{itemize}[leftmargin=1.4em,topsep=2pt]
  \item Simple chat prompts prefer cheap/fast models such as Haiku, mini, flash,
        nano, lite, small, or similar tiers.
  \item Coding, debugging, implementation, research, scraping, planning,
        review, and analysis prompts prefer stronger models such as Sonnet,
        Opus, GPT-4/4o, large/pro/max, 70B+, or similar tiers.
  \item If the exact desired tier is unavailable, routing falls back through
        mid-tier or any available connected model.
  \item If provider settings or model switching hangs, the auto route is skipped
        after a short timeout and task submission continues.
\end{itemize}

\subsection{Local model download and runtime UX}
The Providers panel now places a local model download box directly below the
device checker. It lets the user pick a suggested/cached model, download it,
run a cached model, stop the local server, and choose a runtime target:
\textbf{Auto}, \textbf{GPU}, \textbf{WebGPU}, or \textbf{CPU}. The UI shows
downloading progress, installed state, failed state, and server running/idle
state. Manual GPU index selection (GPU~0/GPU~1) is not supported by the current
runtime surface yet; GPU modes use automatic GPU selection.

\subsection{AI questions}
When the agent sends a question permission request, the execution UI renders a
compact inline \textbf{DigiBull Question} card with the prompt, optional choices,
and a typed-answer field. This makes SRIM's clarification requests visible in
the conversation flow instead of looking like a silent hang.

% ---------------------------------------------------------------
\section{Build and Run}\label{sec:run}
\begin{lstlisting}
# 1. Dependencies (Node 24 recommended)
pnpm install
# 2. Backend daemon (optionally enable inbound integration API)
$env:ACCOMPLISH_ENABLE_INTEGRATION_API='1'   # optional
pnpm -F "@accomplish/daemon" dev             # -> Listening on port 9234
# 3. Web UI
pnpm dev:web                                 # -> http://localhost:5173
\end{lstlisting}
A convenience launcher (\verb|LAUNCH-WEB.bat|) and a manual checklist
(\verb|VERIFY_RESULTS.txt|) are included at the project root.

% ---------------------------------------------------------------
\section{Artifacts}
\begin{itemize}[leftmargin=1.4em,topsep=2pt]
  \item \verb|LOCAL/SRIM-UI/accomplish_Transformed_Chaithanya_Original.zip| ---
        original Chaithanya web transform baseline archive.
  \item \verb|LOCAL/migration_inventory.md| --- Electron/Node API inventory.
  \item \verb|.../accomplish_Transformed_Nr/missing_features_gap_analysis.md| ---
        capability gap analysis (buckets A--D).
  \item \verb|LOCAL/exe_0.4.14_feature_delta.md| --- 0.4.14 vs source delta.
  \item \verb|LOCAL/exe-inspect-0.4.14/| --- extracted asar, file list, bundles.
  \item \verb|LOCAL/docs/MCP_GUIDE.md| --- MCP usage guide.
  \item \verb|LOCAL/docs/DIGIBULL_UI_UPDATE.md| --- DigiBull UI, Auto model,
        local model, and question-card usage notes.
\end{itemize}

\end{document}
